% Advanced Practice Examination: Game Theory and Information Economics
% Questions only. Solutions are in Micro1Q_solutions.tex.
% Compile with:
%   latexmk -pdf -outdir=aux Micro1Q.tex

\documentclass[12pt,a4paper]{article}

% Packages
\usepackage[utf8]{inputenc}
\usepackage[margin=1in]{geometry}
\usepackage{amsmath}
\usepackage{amssymb}
\usepackage{amsfonts}
\usepackage{bm}
\usepackage{mathtools}
\usepackage{booktabs}
\usepackage{tabularx}
\usepackage{enumitem}
\usepackage{hyperref}

% Enumerate style for sub-tasks
\setlist[enumerate,1]{label=\arabic*., leftmargin=*}

\title{\textbf{Advanced Practice Examination:}\\
       \Large Game Theory and Information Economics}
\author{ }
\date{ }

\begin{document}

\maketitle

\section{Simultaneous Games and Dominance Solvability}

The bedrock of non-cooperative game theory lies in the concepts of iterated dominance and Nash Equilibrium. In environments with discrete strategy spaces, the process of Iterated Elimination of Strictly Dominated Strategies (IESDS) allows a researcher to winnow down the set of plausible outcomes by assuming only that players are rational. When supplemented by the elimination of weakly dominated strategies (IEWDS), we can often reach a unique prediction. However, these stable equilibrium points---where no player has a unilateral incentive to deviate---frequently fail to achieve social optimality, illustrating the inherent tension between individual rationality and collective efficiency.

\subsection*{Problem 1.1: The 4x4 Strategic Matrix}

Consider a two-player normal-form game where Player~A (Row) chooses
\(
  s_A \in \{a_1, a_2, a_3, a_4\}
\)
and Player~B (Column) chooses
\(
  s_B \in \{b_1, b_2, b_3, b_4\}
\).
The payoff matrix \((u_A, u_B)\) is defined as follows:

\begin{table}[h!]
  \centering
  \renewcommand{\arraystretch}{1.1}
  \begin{tabular}{lcccc}
    \toprule
    $A \setminus B$ & $b_1$ & $b_2$ & $b_3$ & $b_4$ \\
    \midrule
    $a_1$ & $(4,0)$ & $(2,2)$ & $(8,0)$ & $(2,1)$ \\
    $a_2$ & $(1,0)$ & $(2,1)$ & $(0,6)$ & $(1,1)$ \\
    $a_3$ & $(3,3)$ & $(1,2)$ & $(6,1)$ & $(1,2)$ \\
    $a_4$ & $(5,3)$ & $(0,2)$ & $(1,1)$ & $(2,2)$ \\
    \bottomrule
  \end{tabular}
\end{table}

\noindent\textbf{Tasks.}
\begin{enumerate}
  \item \textbf{Iterated Dominance.} Perform multiple rounds of IESDS. Clearly identify which strategies are eliminated and by which dominating strategy. Subsequently, apply IEWDS to the reduced matrix to find the unique equilibrium prediction.
  \item \textbf{Pareto Efficiency.} Identify all strategy profiles in the original \(4\times 4\) matrix that are Pareto efficient.
  \item \textbf{Analytical Evaluation.} Explain the ``So What?'' layer: Why does the equilibrium reached through dominance solvability result in a payoff profile that is Pareto dominated by other outcomes? Relate this to the ``Fishermen's Dilemma'' context.
\end{enumerate}

\bigskip
\noindent\textbf{Transition.} While discrete choices provide a clear winnowing logic, strategic variables in industrial organization---such as price or quantity---often exist on a continuum, necessitating the use of calculus and reaction functions to find equilibria.

\section{Sequential Interaction, Backward Induction, and Bargaining}

In sequential games, the timing of moves grants players the ability to commit to certain paths, yet the credibility of such commitments is paramount. Subgame Perfect Nash Equilibrium (SPNE) serves as a refinement of the Nash Equilibrium by requiring that players act optimally in every subgame, effectively filtering out non-credible threats that a player would not actually execute if that node were reached.

\subsection*{Problem 2.1: Extensive Form and Forward Induction}

Consider a three-player extensive-form game with \(7\) subgames. Player~1 moves first and can choose to Burn Money (\(B\)) or Not Burn (\(N\)).
\begin{itemize}
  \item If Player~1 chooses \(N\), a subgame follows where Player~1 chooses between \(A\) and \(B\). If \(A\), Player~2 chooses \(C\) or \(D\) (where \(D\) yields \((0,0,0)\)). If \(C\), Player~1 chooses \(U\) or \(V\). If \(V\), Player~3 chooses \(X\) or \(Y\).
  \item If Player~1 chooses \(B\), all of Player~1's terminal payoffs in the subsequent subgame are reduced by \(K = 1\).
  \item The terminal payoffs for the sequence \((N, A, C, U)\) are \((4, 2, 0)\) and for \((N, B, E, W)\) are \((10, -4, 0)\).
\end{itemize}

\subsection*{Problem 2.2: Finite Horizon Bargaining}

Two players must divide a surplus of \(1\) over \(T = 3\) periods. In \(T = 1\), Player~1 proposes \((x, 1-x)\). If Player~2 rejects, the surplus is discounted by \(\delta\) and Player~2 proposes in \(T = 2\). If Player~1 rejects, they move to \(T = 3\) where Player~1 makes the final offer. If rejected at \(T = 3\), both receive \(0\).

\noindent\textbf{Tasks.}
\begin{enumerate}
  \item Using backward induction, solve for the unique SPNE of the \(T = 3\) bargaining game.
  \item \textbf{Analytical Task.} Evaluate the impact of the discount factor \(\delta\) on the first-mover advantage. What is the strategic consequence if \(T \to \infty\)? Show the relationship between the equilibrium split and the formula \(\frac{1}{1+\delta}\).
\end{enumerate}

\bigskip
\noindent\textbf{Transition.} The certainty of sequential moves gives way to strategic unpredictability when players must randomize their actions to prevent being outguessed by opponents.

\section{Mixed Strategies, Security Payoffs, and Zero-Sum Dynamics}

In strictly competitive environments, randomization is not merely an option but a strategic necessity. The Minmax Theorem, established by von~Neumann, demonstrates that in zero-sum games, a player's security payoff (maxmin) is identical to the equilibrium payoff. By randomizing, a player ensures a ``Value'' for the game that cannot be diminished by any action the opponent takes.

\subsection*{Problem 3.1: Military Attack (Strictly Competitive)}

A commander (Player~1) can strike targets \(A\), \(B\), or \(C\). A defender (Player~2) stationing a battery at the same target successfully intercepts the attack. The payoffs (enemy troop losses for Player~1, gains for Player~2) are represented by the following zero-sum matrix:

\begin{table}[h!]
  \centering
  \renewcommand{\arraystretch}{1.1}
  \begin{tabular}{lccc}
    \toprule
    $1 \setminus 2$ & $A$ & $B$ & $C$ \\
    \midrule
    $A$ & $0$ & $4$ & $4$ \\
    $B$ & $3$ & $0$ & $3$ \\
    $C$ & $2$ & $2$ & $0$ \\
    \bottomrule
  \end{tabular}
\end{table}

\noindent\textbf{Tasks.}
\begin{enumerate}
  \item Prove that no pure-strategy Nash Equilibrium exists for this game.
  \item Calculate the Maxmin (security) strategies for Player~1. Show the system of equations derived from making Target \(A\), \(B\), and \(C\)'s expected payoffs equal.
  \item Determine the Value \(v\) of the game.
  \item \textbf{Analytical Task.} Compare the security strategies of this zero-sum game to a non-zero-sum ``Inspection Game.'' Explain why the Minmax Theorem's equality \(v_1 = -v_2\) fails in the latter.
\end{enumerate}

\bigskip
\noindent\textbf{Transition.} Having mastered full-information environments, we now consider markets where one party possesses hidden knowledge about the quality of the good being traded.

\section{Adverse Selection and Market Signaling}

Asymmetric information regarding quality \(\theta\) can lead to the ``Lemons Problem,'' where the inability of buyers to distinguish high from low quality leads to a collapse in trade. If high-quality sellers cannot be adequately compensated, they exit the market, leaving only low-quality types and potentially destroying the gains from trade entirely.

\subsection*{Problem 4.1: The Used Car Market with Disclosure}

Quality \(\theta\) is uniformly distributed on \([0, 2]\). A seller's valuation is
\(
  V_S(\theta) = \theta
\)
and a buyer's valuation is
\(
  V_B(\theta) = \theta + a
\).

\noindent\textbf{Tasks.}
\begin{enumerate}
  \item \textbf{TIOLI Offer.} Solve for the buyer's optimal take-it-or-leave-it price \(p^*\) and identify the threshold of \(a\) that ensures a trade probability of \(1\).
  \item \textbf{Signaling Disclosure.} Suppose the seller can pay a fixed cost \(c\) to reveal \(\theta\) to the buyer. If the seller reveals, the buyer offers the full valuation \(V_B(\theta)\). If the seller does not reveal, they are subject to the buyer's price \(p^*\).
\end{enumerate}

\noindent\textbf{Analytical Task.} Derive the condition on \(c\) such that a seller of quality \(\theta\) will choose to reveal. Evaluate how the ``information rent'' captured by the seller changes with \(a\).

\bigskip
\noindent\textbf{Transition.} Information asymmetries extend beyond hidden types to hidden actions, where incentives must be designed to align the agent's behavior with the principal's goals.

\section{Moral Hazard and Monopolistic Screening}

Moral hazard arises when an agent's effort is unobservable, creating a trade-off between risk-sharing and incentives. In a Principal--Agent framework, the Principal must design a contract that satisfies the agent's Participation Constraint (PC) and Incentive Compatibility (IC) constraint, ensuring the agent chooses the high-effort action despite its private cost.

\subsection*{Problem 5.1: Optimal Monitoring and Agency Costs}

A Principal (Employer) hires an Agent (Worker).
\begin{itemize}
  \item Worker: Can choose to Work \((c = 50, \text{Output } Y = 200)\) or Not Work \((c = 0, Y = 0)\).
  \item Employer: Pays wage \(w = 100\). Can monitor with probability \(m \in [0,1]\) at a cost of \(10m\). If caught shirking, the wage is \(0\).
  \item Reservation Utility: The agent's outside option is \(U_0 = 0\).
\end{itemize}

\noindent\textbf{Tasks.}
\begin{enumerate}
  \item Derive the Incentive Compatibility (IC) and Participation Constraint (PC) for the agent.
  \item Solve for the Principal's optimal choice of \(m\).
  \item \textbf{Efficiency Analysis.} Compare the Principal's profit to the first-best scenario (where effort is observable). Calculate the Agency Cost and explain its origin in the source context.
\end{enumerate}

\bigskip
\noindent\textbf{Transition.} This progression from pure game theory to information-contingent applications underscores the necessity of modeling strategic interaction to understand modern market failures.

\end{document}

